\scanpage{91}
\mspage{87}
% archive sheet number: 90

% Continuation of the block cancelled by two diagonal strokes.
\[
\begin{aligned}
\deleted{\alpha'(x)=\operatorname{supp}^{\circ}_{E'}\alpha(x)}\\
\deleted{\beta'(x)=\operatorname{supp}^{\circ}_E\beta(x).}
\end{aligned}
\]
\deleted{De plus, pour \unsure{$\alpha$} donné $\alpha$ (\inserted{ou $\alpha'$} ou $\alpha_{\unsure{*|\ell}}$ donné), il n'existe \inserted{au plus} qu'un seul $\beta_*$ qui}

\deleted{Si $\beta_1,\beta_2$ sont deux \unsure{transposés} de $\alpha$, on a $\beta'_1=\beta'_2,\ \beta_{1*}=\beta_{2*}$, et \unsure{itou} \unsure{en échangeant|en échange} \inserted{du côté} $\alpha$ si on a deux transposés $\alpha_1,\alpha_2$ de $\beta$.}

\uline{Dém} On a évidemment $(\mathrm{ii})\Rightarrow(\mathrm{i})$, en faisant $A=\{x\},\ A'=\{x'\}$, et $(\mathrm{i})\Rightarrow(\mathrm{ii})$, en notant que
\[
\begin{aligned}
A\mid\circ\mid\bar\beta(A')&\ \Longleftrightarrow\ x\mid\circ\mid\unsure{\overleftarrow\beta|\bar\beta}(x') &&\forall x\in A,\ x'\in A'\\
\unsure{\cdot}\quad\alpha(A)\mid\circ\mid A'&\ \Longleftrightarrow\ \alpha(x)\mid\circ\mid x' &&\rule{2.7cm}{0.3pt}\ .
\end{aligned}
\]

\deleted{Quand \inserted{en plus} $A,A'\in\Sigma_E$, l'équivalence \unsure{des $2$ mbrs de (ii) reste}}
Notons que les deux \unsure{mbrs|membres} de (ii) ne changent pas quand on y remplace l'un ou l'autre des termes dans la relation $\mid\circ\mid$, ou les deux, par son support. Quand $A\in\Sigma_E,A'\in\Sigma_{E'}$, cela montre que cette équivalence \deleted{signifie} équivaut à celle de (iii), donc $(\mathrm{ii})\Rightarrow(\mathrm{iii})\inserted{\text{a)}}$. \inserted{On verra b) plus bas.} \unsure{Enfin}, prouvons que $(\mathrm{iii})\Rightarrow(\mathrm{ii})$ \unsure{(p. ex.)|(p. 86)}. On écrit la relation (ii) à prouver sous la forme équivalente :
